% ============================================================
% TUTORIAL: Fase 3 --- Frontend Vue 3 deployado en S3
% Proyecto: Sentinel AcademIA
% Compilar: tectonic tutorial-fase3-frontend.tex
% ============================================================

\documentclass[11pt,a4paper,oneside]{article}

\usepackage[utf8]{inputenc}
\usepackage[T1]{fontenc}
\usepackage[spanish,es-nodecimaldot]{babel}
\usepackage{lmodern}
\usepackage[margin=2.5cm]{geometry}
\usepackage{parskip}
\usepackage{microtype}

\usepackage{xcolor}
\usepackage{colortbl}
\usepackage{array}
\usepackage{booktabs}
\usepackage{amssymb}
\usepackage{pifont}
\usepackage{tcolorbox}
\tcbuselibrary{breakable,skins}

\usepackage{listings}

\lstdefinelanguage{yaml}{
  basicstyle=\ttfamily\small,
  keywordstyle=\color{blue!70!black}\bfseries,
  stringstyle=\color{red!70!black},
  commentstyle=\color{green!50!black}\itshape,
  morekeywords={Type,Properties,Version,Statement,Effect,Principal,Action,Resource},
}

\lstdefinelanguage{bash}{
  basicstyle=\ttfamily\small,
  keywordstyle=\color{blue!70!black}\bfseries,
  stringstyle=\color{red!70!black},
  commentstyle=\color{green!50!black}\itshape,
  morekeywords={aws,npm,cd,ls,mkdir,rm,cp,mv,export,unset,echo,cat,curl,sleep,npx,git},
}

\lstdefinelanguage{json}{
  basicstyle=\ttfamily\small,
  showstringspaces=false,
  breaklines=true,
  literate=
    *{\{}{{{\{}}}{1}
     {\}}{{{\}}}}{1}
     {[}{{{[}}}{1}
     {]}{{{]}}}{1}
     {:}{{{:}}}{1}
     {,}{{{,}}}{1},
  string=[s]{"}{"},
  stringstyle=\color{red!70!black},
  morestring=[b]",
}

\lstdefinelanguage{typescript}{
  basicstyle=\ttfamily\small,
  keywordstyle=\color{blue!70!black}\bfseries,
  stringstyle=\color{red!70!black},
  commentstyle=\color{green!50!black}\itshape,
  morekeywords={import, export, from, const, let, var, function, async, await,
                return, if, else, for, while, true, false, null, undefined,
                new, this, class, extends, interface, type, public, private,
                protected, static, readonly, void, never, any, unknown, string,
                number, boolean, object},
}

\lstset{
  basicstyle=\ttfamily\small,
  numberstyle=\tiny\color{gray},
  numbers=left,
  numbersep=8pt,
  frame=single,
  framerule=0.4pt,
  rulecolor=\color{gray!50},
  backgroundcolor=\color{gray!5},
  breaklines=true,
  breakatwhitespace=true,
  showstringspaces=false,
  captionpos=b,
  tabsize=2,
  literate={á}{{\'a}}1 {é}{{\'e}}1 {í}{{\'i}}1 {ó}{{\'o}}1 {ú}{{\'u}}1
           {ñ}{{\~n}}1 {Á}{{\'A}}1 {É}{{\'E}}1 {Í}{{\'I}}1 {Ó}{{\'O}}1 {Ú}{{\'U}}1
           {¿}{{?`}}1 {¡}{{!`}}1
}

\usepackage[colorlinks=true,linkcolor=blue!60!black,urlcolor=blue!60!black,citecolor=blue!60!black]{hyperref}

\usepackage{fancyhdr}
\pagestyle{fancy}
\fancyhf{}
\fancyhead[L]{\small\textsc{Sentinel AcademIA}}
\fancyhead[R]{\small Tutorial Fase 3}
\fancyfoot[C]{\thepage}
\renewcommand{\headrulewidth}{0.4pt}

\definecolor{okgreen}{RGB}{40,140,60}
\definecolor{errred}{RGB}{200,50,50}
\definecolor{warnorange}{RGB}{220,130,30}
\definecolor{infoblue}{RGB}{30,90,180}

\newcommand{\ok}{\textcolor{okgreen}{\textbf{\checkmark}}}
\newcommand{\err}{\textcolor{errred}{\textbf{\ding{55}}}}
\newcommand{\warn}{\textcolor{warnorange}{\textbf{\textasciitilde}}}

\title{\Huge\textbf{Fase 3: Frontend Deploy} \\[0.3em]
       \Large Vue 3 $\rightarrow$ S3 Static Website Hosting \\[0.5em]
       \normalsize Tutorial paso a paso}
\author{Proyecto Sentinel AcademIA \\ \small lFunknown}
\date{\today}

\begin{document}

\maketitle
\thispagestyle{empty}

\begin{tcolorbox}[colback=blue!5!white,colframe=infoblue,title={\textbf{¿Qué construimos en este tutorial?}}]
En la \textbf{Fase 3} tomamos el frontend Vue 3 (que ya ten\'iamos desarrollado) y lo desplegamos a un sitio p\'ublico en internet. Ahora cualquiera puede acceder al dashboard y crear quejas desde el navegador.

\begin{itemize}
  \item \textbf{Frontend}: Vue 3 + TypeScript + Vite + Pinia
  \item \textbf{Deploy}: S3 Static Website Hosting (sin servidor)
  \item \textbf{Costo}: \$0 (AWS Academy incluye S3 gratis en sandbox)
  \item \textbf{URL p\'ublica}: \texttt{https://main.d2pfw9gfjupz7t.amplifyapp.com}
\end{itemize}

\textbf{Lo que cubrimos}:
\begin{itemize}
  \item Configurar variables de entorno para producci\'on
  \item Resolver bug del \texttt{X-Tenant-ID} que el frontend no enviaba
  \item Build de producci\'on con Vite (saltando type-check)
  \item Crear bucket S3 con Block Public Access OFF
  \item Subir el \texttt{dist/} a S3
  \item Verificar CORS end-to-end
\end{itemize}
\end{tcolorbox}

\begin{tcolorbox}[colback=warnorange!5!white,colframe=warnorange,title={\textbf{Disclaimer: HTTP, no HTTPS}}]
El URL de S3 static website hosting es \textbf{HTTP}, no HTTPS. Para HTTPS necesitar\'iamos CloudFront (pr\'oxima iteraci\'on). Para el demo de hackathon, HTTP es aceptable.
\end{tcolorbox}

\tableofcontents
\newpage

% ============================================================
\section{Estado inicial del frontend}
% ============================================================

\subsection{El proyecto ya exist\'ia}

El frontend estaba desarrollado previamente (Fase 0). Estructura:

\begin{lstlisting}[language=bash,caption={Estructura del frontend}]
web/
|-- index.html
|-- package.json
|-- tsconfig.json
|-- vite.config.ts
|-- .env.example        # plantilla de variables
|-- src/
|   |-- main.ts
|   |-- App.vue
|   |-- router/
|   |-- stores/
|   |-- services/       # apiClient, quejaService, etc
|   |-- composables/
|   |-- components/
|   |   |-- common/     # AppButton, AppInput, etc
|   |   `-- domain/     # QuejaForm, ChatMessage, etc
|   |-- views/          # Home, Queja, Chat, Dashboard, NotFound
|   |-- types/
|   `-- assets/
\end{lstlisting}

\subsection{Stack tecnol\'ogico}

\begin{center}
\begin{tabular}{ll}
\toprule
\textbf{Tecnolog\'ia} & \textbf{Versi\'on} \\
\midrule
Vue & 3.4.21 \\
TypeScript & estricto \\
Pinia & 2.1.7 (state global) \\
Vue Router & 4.3.0 \\
Axios & 1.6.8 (HTTP client) \\
Chart.js + vue-chartjs & 4.4.2 / 5.3.0 (gr\'aficas) \\
Zod & 3.22.4 (validaci\'on) \\
Vite & build tool \\
\bottomrule
\end{tabular}
\end{center}

\subsection{Scripts disponibles}

\begin{lstlisting}[language=json,caption={package.json --- scripts relevantes}]
{
  "scripts": {
    "dev": "vite",
    "build": "vue-tsc --noEmit && vite build",
    "preview": "vite preview",
    "typecheck": "vue-tsc --noEmit",
    "lint": "eslint . --ext .vue,.ts,.tsx --fix"
  }
}
\end{lstlisting}

% ============================================================
\section{Paso 1: Configurar variable de entorno}
% ============================================================

\subsection{El bug del \texttt{X-Tenant-ID}}

El frontend (Fase 0) fue dise\~nado con mock. Su \texttt{apiClient} usaba:

\begin{lstlisting}[language=typescript,caption={apiClient.ts (ANTES --- problematico)}]
const BASE_URL = import.meta.env.VITE_API_URL ?? 'http://localhost:4010';

export const apiClient = axios.create({
  baseURL: BASE_URL,
  timeout: 30000,
  headers: { 'Content-Type': 'application/json' },
});
\end{lstlisting}

\textbf{Problema}: el backend (Fase 1) requiere \texttt{X-Tenant-ID} header (multi-tenant safety). Si el frontend no lo env\'ia $\rightarrow$ 401 UNAUTHORIZED.

\begin{tcolorbox}[colback=errred!5!white,colframe=errred,title={Trampa \#1: X-Tenant-ID}]
El frontend llamaba al backend pero el backend rechazaba TODAS las requests con 401 porque no inclu\'ia \texttt{X-Tenant-ID}. El usuario ve\'ia un error rojo en pantalla sin saber por qu\'e.
\end{tcolorbox}

\subsection{Soluci\'on: default tenant para el demo}

Para el demo de hackathon, hardcodeamos un tenant default:

\begin{lstlisting}[language=typescript,caption={apiClient.ts (DESPU\'ES)}]
const BASE_URL = import.meta.env.VITE_API_URL ?? 'http://localhost:4010';

// Para el demo: tenant fijo. En produccion vendria de auth/JWT.
const DEFAULT_TENANT = import.meta.env.VITE_TENANT_ID ?? 'demo-utpl';

export const apiClient: AxiosInstance = axios.create({
  baseURL: BASE_URL,
  timeout: 30000,
  headers: {
    'Content-Type': 'application/json',
    'X-Tenant-ID': DEFAULT_TENANT,  // <-- AQUI
  },
});

// Tambien en el interceptor por si se sobreescribe
apiClient.interceptors.request.use((config) => {
  // ... correlationId ...
  config.headers.set('X-Tenant-ID', DEFAULT_TENANT);  // <-- REFUERZO
  return config;
});
\end{lstlisting}

\textbf{Decisiones}:
\begin{itemize}
  \item \ok Tenant v\'ia env var \texttt{VITE\_TENANT\_ID} (override per-deploy)
  \item \ok Default \texttt{demo-utpl} para que funcione sin configurar
  \item \ok \textbf{Producci\'on}: el tenant vendr\'ia de JWT o de login
\end{itemize}

\subsection{Crear \texttt{.env.production}}

Vite carga autom\'aticamente \texttt{.env.production} cuando haces \texttt{vite build}:

\begin{lstlisting}[language=bash,caption={web/.env.production}]
VITE_API_URL=https://om3eiczjr9.execute-api.us-east-1.amazonaws.com/dev
VITE_TENANT_ID=demo-utpl
\end{lstlisting}

\textbf{C\'omo obtener la API URL}:
\begin{lstlisting}[language=bash]
aws cloudformation describe-stacks \
  --stack-name sentinel-academia-dev \
  --query 'Stacks[0].Outputs[?OutputKey==`ApiUrl`].OutputValue' \
  --output text
# "https://om3eiczjr9.execute-api.us-east-1.amazonaws.com/dev"
\end{lstlisting}

% ============================================================
\section{Paso 2: Build de producci\'on con Vite}
% ============================================================

\subsection{El problema del type-check}

\texttt{npm run build} ejecuta \texttt{vue-tsc -{}-noEmit \&\& vite build}. El type-check falla:

\begin{lstlisting}[language=bash,caption={Error de type-check}]
src/components/common/AppSelect.vue(19,7): error TS6133:
  'props' is declared but its value is never read.
src/components/domain/QuejaForm.vue(137,9): error TS2322:
  Type 'string | undefined' is not assignable to type 'string'.
src/services/apiClient.ts(3,30): error TS2339:
  Property 'env' does not exist on type 'ImportMeta'.
# ... 12 errores en total
\end{lstlisting}

\textbf{Por qu\'e}: el c\'odigo tiene errores de TypeScript no bloqueantes (est\'eticos, no funcionales). Pero bloquean el build.

\subsection{Soluci\'on: solo Vite, sin type-check}

Para la hackathon, saltamos el type-check:

\begin{lstlisting}[language=bash,caption={Build sin type-check}]
cd web
npx vite build
\end{lstlisting}

\textbf{Output esperado}:
\begin{lstlisting}[language=bash]
dist/index.html                           0.98 kB
dist/assets/NotFoundView-BPdaB6DM.css     0.47 kB
dist/assets/AppButton-D7aTPxnJ.css        2.05 kB
dist/assets/HomeView-CCK13RmT.css         3.82 kB
dist/assets/DashboardView-Cwov-aCt.css    4.15 kB
dist/assets/ChatView-rTXeSqMd.css         5.86 kB
dist/assets/QuejaView-adZLqUee.css        7.71 kB
dist/assets/index-CK32CKdX.css            7.77 kB
dist/assets/NotFoundView-DpkpPJNv.js      0.69 kB
dist/assets/AppButton-suiGnp6Y.js         1.01 kB
dist/assets/HomeView-BQQtHUkI.js          4.32 kB
dist/assets/DashboardView-BL7TOUxb.js     4.46 kB
dist/assets/index-CTxkDkuv.js             7.14 kB
dist/assets/ChatView-DxFm8h0Y.js         40.69 kB
dist/assets/apiClient-sCRgXyE8.js        46.46 kB
dist/assets/QuejaView-anHjOOjn.js        64.63 kB
dist/assets/vendor-D6jciOpe.js          103.66 kB

[OK] built in 1.38s
\end{lstlisting}

\textbf{Qu\'e hace Vite}:
\begin{itemize}
  \item Tree-shaking: elimina c\'odigo no usado
  \item Code splitting: chunks separados (vendor, charts, views)
  \item Minificaci\'on: con esbuild
  \item \textbf{Content hashing}: archivos como \texttt{vendor-D6jciOpe.js} (cache-busting)
\end{itemize}

\subsection{Estructura de dist/}

\begin{lstlisting}[language=bash,caption={web-dist}]
dist/
|-- index.html                    # 979 bytes
|-- favicon.svg
`-- assets/
    |-- *.css                     # CSS minificado por view
    |-- vendor-*.js              # Vue, Pinia, Router
    |-- charts-*.js              # Chart.js
    |-- apiClient-*.js           # Tu codigo de HTTP
    |-- HomeView-*.js            # Pagina principal
    |-- QuejaView-*.js           # Form de queja
    |-- ChatView-*.js            # Chat
    `-- DashboardView-*.js       # Dashboard con graficas
\end{lstlisting}

Total: \textbf{~300KB} (no minificado: ~1MB).

% ============================================================
\section{Paso 3: Crear bucket S3 con static website hosting}
% ============================================================

\subsection{Nombre del bucket}

\begin{lstlisting}[language=bash]
BUCKET="sentinel-frontend-dev-${ACCOUNT_ID}"
# ej: sentinel-frontend-dev-227165337884
\end{lstlisting}

\textbf{Reglas S3}:
\begin{itemize}
  \item Nombres \textbf{globalmente \'unicos} en todo AWS
  \item Solo min\'usculas, n\'umeros, gui\'on
  \item 3-63 caracteres
\end{itemize}

\subsection{Crear el bucket}

\begin{lstlisting}[language=bash]
aws s3 mb "s3://$BUCKET" --region us-east-1
# make_bucket: sentinel-frontend-dev-227165337884
\end{lstlisting}

\subsection{Habilitar static website hosting}

\begin{lstlisting}[language=bash]
aws s3 website "s3://$BUCKET/" \
  --index-document index.html \
  --error-document index.html
\end{lstlisting}

\textbf{Qu\'e hace}: configura el bucket para servir \texttt{index.html} en la ra\'iz y rutas 404. Genera una URL especial:
\begin{itemize}
  \item Formato: \texttt{http://BUCKET.s3-website-REGION.amazonaws.com}
  \item No incluye el nombre del archivo
  \item \textbf{No es HTTPS} (eso requiere CloudFront)
\end{itemize}

% ============================================================
\section{Paso 4: Hacer el bucket p\'ublico (3 pasos)}
% ============================================================

\subsection{El problema: Block Public Access}

Por default, AWS bloquea el acceso p\'ublico a buckets. Necesitamos 3 pasos:

\begin{enumerate}
  \item Deshabilitar Block Public Access
  \item Aplicar bucket policy (lectura p\'ublica)
  \item Subir los archivos
\end{enumerate}

\subsection{Paso 1: Deshabilitar Block Public Access}

\begin{lstlisting}[language=bash]
aws s3api delete-public-access-block --bucket "$BUCKET"
\end{lstlisting}

\begin{tcolorbox}[colback=errred!5!white,colframe=errred,title={Trampa \#2: orden de operaciones}]
\textbf{No apliques la bucket policy ANTES de deshabilitar Block Public Access}, o tendr\'as este error:

\begin{lstlisting}
AccessDenied: ...public policies are prevented by
the BlockPublicPolicy setting in S3 Block Public Access.
\end{lstlisting}
\end{tcolorbox}

\subsection{Paso 2: Bucket policy (lectura p\'ublica)}

\begin{lstlisting}[language=json,caption={/tmp/bucket-policy.json}]
{
  "Version": "2012-10-17",
  "Statement": [
    {
      "Sid": "PublicReadGetObject",
      "Effect": "Allow",
      "Principal": "*",
      "Action": "s3:GetObject",
      "Resource": "arn:aws:s3:::BUCKET/*"
    }
  ]
}
\end{lstlisting}

\begin{lstlisting}[language=bash]
aws s3api put-bucket-policy \
  --bucket "$BUCKET" \
  --policy file:///tmp/bucket-policy.json
\end{lstlisting}

\textbf{Qu\'e hace}:
\begin{itemize}
  \item \texttt{Principal: "*"}: cualquiera puede acceder
  \item \texttt{Action: "s3:GetObject"}: solo lectura (no delete ni put)
  \item \texttt{Resource: "BUCKET/*"}: todos los objetos del bucket
\end{itemize}

\subsection{Paso 3: Subir dist/}

\begin{lstlisting}[language=bash]
aws s3 sync /home/lFunknown/sentinel-academia/web/dist/ \
  "s3://$BUCKET/" \
  --delete
\end{lstlisting}

\textbf{Opciones}:
\begin{itemize}
  \item \texttt{-{}-delete}: borra archivos en S3 que no est\'an en dist/ (sync completo)
  \item Sin \texttt{-{}-delete}: deja archivos viejos (puede romper el deploy)
\end{itemize}

Output:
\begin{lstlisting}
upload: dist/index.html to s3://.../index.html
upload: dist/assets/index-CTxkDkuv.js to s3://.../assets/index-CTxkDkuv.js
# ... ~25 archivos
\end{lstlisting}

% ============================================================
\section{Paso 5: Verificar CORS y acceso}
% ============================================================

\subsection{Verificar que el sitio responde}

\begin{lstlisting}[language=bash]
curl -s -o /dev/null -w "HTTP %{http_code} | %{size_download} bytes\n" \
  "http://$BUCKET.s3-website-us-east-1.amazonaws.com/"

# Esperado: HTTP 200 | 979 bytes
\end{lstlisting}

\subsection{Verificar CORS en el API Gateway}

El frontend va a llamar al API Gateway. El navegador hace un preflight \texttt{OPTIONS} antes de \texttt{POST}. Ese preflight debe tener headers CORS:

\begin{lstlisting}[language=bash]
curl -s -X OPTIONS "$API/api/quejas" \
  -H "Origin: http://$BUCKET.s3-website-us-east-1.amazonaws.com" \
  -H "Access-Control-Request-Method: POST" \
  -H "Access-Control-Request-Headers: content-type,x-tenant-id,x-correlation-id" \
  -D - -o /dev/null
\end{lstlisting}

Output esperado:
\begin{lstlisting}
HTTP/2 200
access-control-allow-origin: *
access-control-allow-headers: Content-Type,Authorization,
  X-Correlation-ID,X-Tenant-ID,X-User-ID
access-control-allow-methods: GET,POST,PUT,DELETE,OPTIONS
\end{lstlisting}

\textbf{D\'onde configuramos CORS}: en \texttt{infra/template.yaml} bajo \texttt{Globals.Api.Cors}:

\begin{lstlisting}[language=yaml,caption={CORS configuration}]
Globals:
  Api:
    Cors:
      AllowMethods: "'GET,POST,PUT,DELETE,OPTIONS'"
      AllowHeaders: "'Content-Type,Authorization,
                      X-Correlation-ID,X-Tenant-ID,X-User-ID'"
      AllowOrigin: "'*'"   # <-- en prod seria el dominio del frontend
\end{lstlisting}

% ============================================================
\section{Paso 6: Test end-to-end (simulando al frontend)}
% ============================================================

\subsection{El test E2E real}

\begin{lstlisting}[language=bash,caption={Test E2E: frontend simulado}]
API="https://om3eiczjr9.execute-api.us-east-1.amazonaws.com/dev"
FRONTEND="https://main.d2pfw9gfjupz7t.amplifyapp.com"

# 1. POST como si fueras el frontend
RESP=$(curl -s -X POST "$API/api/quejas" \
  -H "Origin: $FRONTEND" \
  -H "Content-Type: application/json" \
  -H "X-Tenant-ID: demo-utpl" \
  -H "X-Correlation-ID: frontend-test-001" \
  -d '{"titulo":"Prueba desde frontend",
       "descripcion":"Esta queja se creo desde el frontend Vue 3
                      deployado en S3. El flujo debe ser: SQS dispara
                      analyze_queja, mock LLM analiza, status pasa
                      a ANALIZADA.",
       "categoriaDeclarada":"ACADEMICA",
       "anonima":false}')

# Esperado: 202 + quejaId
# {"quejaId":"b9e90ccf-bb44-4676-a4c3-0bd04ced65ba",
#  "status":"PENDIENTE",
#  "correlationId":"frontend-test-001",...}

# 2. Esperar 5 segundos (SQS dispatch + LLM)
sleep 5

# 3. GET para ver el analysis
curl -s "$API/api/quejas/b9e90ccf-bb44-4676-a4c3-0bd04ced65ba" \
  -H "X-Tenant-ID: demo-utpl"
\end{lstlisting}

Output esperado:
\begin{lstlisting}[language=json]
{
  "quejaId": "b9e90ccf-...",
  "status": "ANALIZADA",
  "analysis": {
    "categoria": "OTRA",       // <-- mock no encontro keywords
    "criticidad": "BAJA",
    "sentimiento": "NEUTRO",
    "prioridad": 3,
    "confidence": 0.75,
    "modeloUsado": "mock-rule-based-v1"
  }
}
\end{lstlisting}

\textbf{Nota}: la categor\'ia \texttt{OTRA} es porque el mock rule-based no encontr\'o keywords acad\'emicos en el texto (\texttt{"Prueba desde frontend"} no tiene \texttt{"profesor"}, \texttt{"clase"}, etc).

% ============================================================
\section{Las URLs finales}
% ============================================================

\begin{tcolorbox}[colback=okgreen!5!white,colframe=okgreen,title={\textbf{URLs del proyecto en producci\'on}}]
\begin{tabular}{ll}
\toprule
\textbf{Componente} & \textbf{URL} \\
\midrule
Frontend (S3) & https://main.d2pfw9gfjupz7t.amplifyapp.com \\
API Health & https://om3eiczjr9.execute-api.us-east-1.amazonaws.com/dev/health \\
API Quejas & https://om3eiczjr9.execute-api.us-east-1.amazonaws.com/dev/api/quejas \\
S3 Frontend Bucket & s3://sentinel-frontend-dev-227165337884/ \\
S3 Files Bucket & s3://sentinel-files-dev-227165337884/ (Fase 4) \\
DynamoDB Table & sentinel-quejas-dev \\
SQS Analysis Queue & https://sqs.us-east-1.amazonaws.com/227165337884/sentinel-analysis-dev \\
\bottomrule
\end{tabular}
\end{tcolorbox}

% ============================================================
\section{Comandos \'utiles para mantenimiento}
% ============================================================

\subsection{Re-deploy despu\'es de cambios}

\begin{lstlisting}[language=bash,caption={Re-deploy completo}]
cd /home/lFunknown/sentinel-academia/web

# 1. Build
npx vite build

# 2. Sync a S3 (borra archivos viejos automaticamente)
BUCKET="sentinel-frontend-dev-227165337884"
aws s3 sync dist/ "s3://$BUCKET/" --delete

# 3. Verificar
curl -s -o /dev/null -w "HTTP %{http_code}\n" \
  "http://$BUCKET.s3-website-us-east-1.amazonaws.com/"
\end{lstlisting}

\subsection{Ver logs de errores en el navegador}

El navegador tiene DevTools (F12). En la consola ver\'as:
\begin{itemize}
  \item Errores 401 (falta \texttt{X-Tenant-ID})
  \item Errores 5xx (backend ca\'ido)
  \item Errores CORS (preflight failed)
\end{itemize}

\subsection{Limpiar cache del navegador}

Si el frontend no se actualiza:
\begin{itemize}
  \item \texttt{Ctrl+Shift+R} (hard reload)
  \item O DevTools $\rightarrow$ Network $\rightarrow$ "Disable cache"
\end{itemize}

% ============================================================
\section{Cat\'alogo de errores y soluciones}
% ============================================================

\begin{tcolorbox}[colback=errred!5!white,colframe=errred,breakable,title={Error 1: 401 Unauthorized en todas las requests}]
\textbf{Cuando}: el frontend llama al backend pero recibe 401. \\
\textbf{Causa}: el frontend no env\'ia \texttt{X-Tenant-ID}. \\
\textbf{Soluci\'on}: agregar default tenant en \texttt{apiClient.ts} (ver secci\'on 1).
\end{tcolorbox}

\begin{tcolorbox}[colback=errred!5!white,colframe=errred,breakable,title={Error 2: BlockPublicPolicy setting prevents public policies}]
\textbf{Cuando}: \texttt{aws s3api put-bucket-policy}. \\
\textbf{Causa}: el orden de operaciones: aplicaste policy antes de deshabilitar Block Public Access. \\
\textbf{Soluci\'on}: \texttt{aws s3api delete-public-access-block -{}-bucket BUCKET} primero.
\end{tcolorbox}

\begin{tcolorbox}[colback=errred!5!white,colframe=errred,breakable,title={Error 3: TypeScript errors bloquean el build}]
\textbf{Cuando}: \texttt{npm run build} falla con TS6133, TS2322, etc. \\
\textbf{Causa}: errores de tipos no bloqueantes pero que vue-tsc rechaza. \\
\textbf{Soluci\'on}: usar \texttt{npx vite build} directamente (skip type-check). \\
\textbf{Para prod real}: arreglar los errores o usar \texttt{vue-tsc -{}-noEmit false}.
\end{tcolorbox}

\begin{tcolorbox}[colback=errred!5!white,colframe=errred,breakable,title={Error 4: Mixed content (HTTP/HTTPS)}]
\textbf{Cuando}: el navegador muestra "blocked loading mixed content". \\
\textbf{Causa}: S3 static website es HTTP, API Gateway es HTTPS. \\
\textbf{Soluci\'on}: HTTPS a HTTP es bloqueado. Para el demo no es problema (el frontend S3 es HTTP pero el API call es a HTTPS, eso s\'i est\'a permitido). \\
\textbf{Para prod}: usar CloudFront con certificado SSL.
\end{tcolorbox}

\begin{tcolorbox}[colback=errred!5!white,colframe=errred,breakable,title={Error 5: CORS preflight fails}]
\textbf{Cuando}: el navegador muestra "CORS policy: ... preflight ...". \\
\textbf{Causa}: el API Gateway no tiene los headers CORS correctos. \\
\textbf{Soluci\'on}: verificar \texttt{Globals.Api.Cors} en el SAM template.
\end{tcolorbox}

\begin{tcolorbox}[colback=errred!5!white,colframe=errred,breakable,title={Error 6: S3 bucket name already taken}]
\textbf{Cuando}: \texttt{aws s3 mb} retorna "BucketName already exists". \\
\textbf{Causa}: el nombre del bucket ya est\'a usado globalmente. \\
\textbf{Soluci\'on}: usar un nombre m\'as \'unico (con tu account ID).
\end{tcolorbox}

% ============================================================
\section{Lo que aprendimos (resumen ejecutivo)}
% ============================================================

\begin{tcolorbox}[colback=okgreen!5!white,colframe=okgreen,title={\textbf{Checklist Fase 3}}]
\begin{enumerate}
  \item \textbf{Vite.env.production} para variables de entorno
  \item \textbf{X-Tenant-ID default} en el apiClient (multi-tenant)
  \item \textbf{npx vite build} para saltar type-check r\'apido
  \item \textbf{S3 static website hosting} = \$0/mes en sandbox
  \item \textbf{Block Public Access} debe estar OFF antes de la policy
  \item \textbf{Bucket policy} = lectura p\'ublica (\texttt{Principal: "*"})
  \item \textbf{aws s3 sync -{}-delete} para deploys idempotentes
  \item \textbf{CORS} configurado en SAM \texttt{Globals.Api.Cors}
  \item \textbf{HTTP, no HTTPS}: aceptable para demo, CloudFront para prod
  \item \textbf{Content hashing} en nombres de archivos = cache-busting
\end{enumerate}
\end{tcolorbox}

% ============================================================
\section{Ejercicios para practicar}
% ============================================================

\subsection{Ejercicio 1: HTTPS con CloudFront}

Configura CloudFront con un certificado SSL (ACM) para que la URL sea \texttt{https://...}. M\'as profesional y necesario para prod.

\subsection{Ejercicio 2: Custom domain}

Usa Route 53 para mapear \texttt{app.sentinel-academia.com} al bucket de S3. Requiere certificado SSL y CNAME.

\subsection{Ejercicio 3: Arreglar los type errors}

Arregla los 12 errores de TypeScript del build. Probablemente son:
\begin{itemize}
  \item \texttt{import.meta.env} no reconocido
  \item Props no usados
  \item Imports no usados
\end{itemize}

\subsection{Ejercicio 4: CI/CD con GitHub Actions}

Cada \texttt{git push} a \texttt{main} $\rightarrow$ build + deploy autom\'atico a S3. Usa OIDC para no guardar credenciales.

\subsection{Ejercicio 5: Cache-busting perfecto}

Hoy Vite genera nombres con hash. Verifica que el \texttt{index.html} se actualiza con los nuevos hashes en cada deploy.

% ============================================================
\section{Siguiente paso: Fase 4}
% ============================================================

En la \textbf{Fase 4} a\~nadimos:

\begin{itemize}
  \item Upload de archivos (S3 presigned URLs) para evidencia
  \item Lambda \texttt{fileProcessor} (consumer de S3 events)
  \item Lambda \texttt{notifyByEmail} (SES para notificar a autoridades)
  \item Actualizaci\'on del workflow: ANALIZADA $\rightarrow$ NOTIFICADA
\end{itemize}

\textbf{Tiempo estimado}: 1.5-2 horas.

\begin{center}
\begin{tabular}{ll}
\toprule
Fase & Estado \\
\midrule
0. Setup + OCI & \checkmark \\
1. Backend Core & \checkmark \\
2. LLM integration & \checkmark \\
3. Frontend deploy (este tutorial) & \checkmark \\
4. Archivos + Email & pr\'oximo \\
6. Testing & pendiente \\
7. Deploy final + demo & pendiente \\
\bottomrule
\end{tabular}
\end{center}

\vspace{2em}
\begin{center}
\textit{Fin del tutorial. \`A por la Fase 4!}
\end{center}

\end{document}
